% part: intuitionistic-logic
% chapter: semantics
% section: topological-semantics

\documentclass[../../../include/open-logic-section]{subfiles}

\begin{document}

\olfileid{int}{sem}{top}

\olsection{Semantik Topologis}

Cara liyane kanggo menehi semantik logika intuisionistik yaiku
nggunakake konsep matematika topologi.

\begin{defn}
  Upamane $X$ iku himpunan. \emph{Topologi ing~$X$} iku himpunan
  $\Top{O} \subseteq \Pow{X}$ kang nyukupi sipat-sipat ing ngisor iki.
  !!{element}s saka~$\Top{O}$ diarani \emph{himpunan kabuka} saka
  topologi iku. Himpunan $X$ bebarengan karo $\Top{O}$ diarani
  \emph{ruang topologis}.
  \begin{enumerate}
  \item Himpunan kosong lan kabeh ruang iku kabuka: $\emptyset$, $X \in
    \Top{O}$.
  \item Himpunan kabuka katutup tumrap irisan winates: yen $U$, $V \in
    \Top{O}$, mula $U \cap V \in \Top{O}$
  \item Himpunan kabuka katutup tumrap gabungan sembarang: yen $U_i \in
    \Top{O}$ kanggo kabeh $i \in I$, mula $\bigcup \Setabs{U_i}{i \in
      I} \in \Top{O}$.
  \end{enumerate}
\end{defn}

Kita bisa nulis $X$ kanggo topologi yen kumpulan himpunan kabukane
bisa dingerteni saka konteks. Nanging, mung sawisé $X$ diwenehi
himpunan kabuka iku, bisa diarani topologi.

\begin{defn}
  \emph{Model topologis} kanggo logika proposisional intuisionistik
  iku tripel $\mModel{X} = \tuple{X, \Top{O}, V}$ kanthi $\Top{O}$
  topologi ing~$X$ lan $V$ fungsi kang menehi himpunan kabuka ing
  $\Top{O}$ marang saben variabel proposisional.

  Yen diwenehi model topologis~$\mModel{X}$, kita bisa nemtokake
  $\Prop{X}{!A}$ sacara induktif kaya mangkene:
  \begin{enumerate}
  \item $\Prop{X}{\lfalse} = \emptyset$
  \item $\Prop{X}{p} = V(p)$
  \item $\Prop{X}{!A \land !B} = \Prop{X}{!A} \cap \Prop{X}{!B}$
  \item $\Prop{X}{!A \lor !B} = \Prop{X}{!A} \cup \Prop{X}{!B}$
  \item $\Prop{X}{!A \lif !B} = \Interior{(X \setminus \Prop{X}{!A}) \cup
    \Prop{X}{!B}}$
  \end{enumerate}
  Ing kene, $\Interior{V}$ nuduhake asil kanggo himpunan
  $V \subseteq X$, yaiku \emph{interior}, gabungan kabeh himpunan kabuka kang
  kakandhut ing kono. Kanthi tembung liya,
  \[
  \Interior{V} = \bigcup \Setabs{U}{U \subseteq V \text{ lan } U \in \Top{O}}.
  \]
\end{defn}

Interior saka himpunan apa wae tansah kabuka, amarga iku gabungan
himpunan kabuka. Mula $\Prop{X}{!A}$ tansah himpunan kabuka.

Sanajan semantik topologis banget abstrak, ana cara kanggo
mikirake kang bisa menehi alasan kanggo semantik iki. Upamane
!!{element}s, utawa ``titik,'' saka $X$ iku titik kang bisa digunakake
kanggo ngevaluasi pranyatan. Himpunan kabeh titik kang $!A$ bener
iku proposisi kang diwedharake dening~$!A$. Ora saben himpunan titik
bisa dadi proposisi; mung !!{element}s saka $\Top{O}$ kang bisa.
$!A \Entails !B$ yen lan mung yen $!B$ bener ing saben titik kang~$!A$
bener, yaiku $\Prop{X}{!A} \subseteq \Prop{X}{!B}$, kanggo kabeh~$X$.
Pranyatan absurd~$\lfalse$ ora tau bener, mula
$\Prop{X}{\lfalse} = \emptyset$.

Kepiye sesambungane proposisi kang diwedharake dening $!B \land !C$,
$!B \lor !C$, lan $!B \lif !C$ karo kang diwedharake dening $!B$
lan~$!C$ supaya hukum-hukum intuisionistik kang valid lumaku?
Tegese, supaya $!A \Proves !B$ yen lan mung yen
$\Prop{X}{!A} \subseteq \Prop{X}{!B}$ kanggo kabeh model topologis.
Kita mbutuhake $\lfalse \Proves !A$ kanggo sembarang $!A$, kang
kasembadan amarga $\emptyset \subseteq U$ kanggo kabeh $U$.
Amarga $!B \land !C \Proves !B$, kita mbutuhake
$\Prop{X}{!B \land !C} \subseteq \Prop{X}{!B}$, lan uga
$\Prop{X}{!B \land !C} \subseteq \Prop{X}{!C}$. Himpunan paling gedhe
kang nyukupi $W \subseteq U$ lan $W \subseteq V$ yaiku $U \cap V$.
Kosokbaline, $!B \Proves !B \lor !C$ lan $!C \Proves !B \lor !C$,
mula kita mbutuhake $\Prop{X}{!B} \subseteq \Prop{X}{!B \lor !C}$ lan
$\Prop{X}{!C} \subseteq \Prop{X}{!B \lor !C}$. Himpunan paling cilik~$W$
kanthi $U \subseteq W$ lan $V \subseteq W$ yaiku $U \cup V$.

Definisi kanggo $\lif$ rada angel: $!A \lif !B$ ngandharake
proposisi paling ringkih kang, yen digabungake karo $!A$, ngentail $!B$.
Yen $!A \lif !B$ digabungake karo $!A$ ngentail~$!B$, iki cetha saka
$(!A \lif !B) \land !A \Proves !B $. Dadi $\Prop{X}{!A \lif !B}$
kudu dadi himpunan kabuka paling gedhe kang nyukupi
$\Prop{X}{!A \lif !B} \cap \Prop{X}{!A} \subseteq \Prop{X}{!B}$,
kang ndadekake definisi ing ndhuwur.

\end{document}
